
% BEGIN SOURCE_BLOCK preamble
\ifdefined\XeTeXversion
\else
\pdfoutput=1
\fi
\documentclass[11pt,a4paper]{article}
\usepackage[T1]{fontenc}
\usepackage{lmodern}
\usepackage{amsmath,amssymb,amsthm,mathtools}
\usepackage[textwidth=440pt,top=27mm,bottom=27mm]{geometry}
\usepackage{microtype}
\usepackage{booktabs,array}
\usepackage{xcolor}
\usepackage{needspace}
\usepackage[colorlinks=true,linkcolor=blue!45!black,citecolor=blue!45!black,urlcolor=blue!45!black]{hyperref}
\hypersetup{pdftitle={Polynomial families of incident flags and explicit off-diagonal Ramsey graphs},pdfauthor={Brecht Verbeken}}
\newcommand{\F}{\mathbb F}
\newcommand{\A}{\mathbb A}
\newcommand{\PP}{\mathbb P}
\newcommand{\Ccal}{\mathcal C}
\newcommand{\Ical}{\mathcal I}
\newcommand{\Wcal}{\mathcal W}
\DeclareMathOperator{\cdeg}{cdeg}
\DeclareMathOperator{\rank}{rank}
\DeclareMathOperator{\Tr}{Tr}
\DeclareMathOperator{\spanop}{span}
\newtheorem{theorem}{Theorem}[section]
\newtheorem{proposition}[theorem]{Proposition}
\newtheorem{lemma}[theorem]{Lemma}
\newtheorem{corollary}[theorem]{Corollary}
\theoremstyle{definition}
\newtheorem{definition}[theorem]{Definition}
\theoremstyle{remark}
\newtheorem{remark}[theorem]{Remark}
\numberwithin{equation}{section}
\setlength{\emergencystretch}{1em}
\begin{document}
\title{Polynomial families of incident flags\\and explicit off-diagonal Ramsey graphs}
\author{Brecht Verbeken\textsuperscript{1,2}\\[0.6em]
\parbox{0.96\textwidth}{\centering\small
\textsuperscript{1}Department of Business Technology and Operations,\\
Data Analytics Laboratory, Vrije Universiteit Brussel (VUB),\\
Pleinlaan 2, 1050 Brussels, Belgium.\\[0.35em]
\textsuperscript{2}imec-SMIT, Vrije Universiteit Brussel,\\
Pleinlaan 9, 1050 Brussels, Belgium.\\[0.45em]
\href{mailto:brecht.verbeken@vub.be}{brecht.verbeken@vub.be}}}
\date{}
\maketitle
% END SOURCE_BLOCK preamble

% BEGIN SOURCE_BLOCK revised_intro
\begin{abstract}
We construct explicit off-diagonal Ramsey graphs from polynomial families of incident point--hyperplane flags. Universal interpolation translates the exact dimension of ordered clique configurations into a bound on their coefficient--label incidence. At the critical dimension, a polynomial separates the forbidden-pair image from its diagonal; deleting the corresponding edges preserves a triangular rank certificate. Restriction of scalars realizes the rational endpoint without rounding loss. The resulting fixed-$s$ exponent has leading scale $s/(2\log_2s)$, and short Frobenius relations give the explicit example $R(16,t)\ge\Omega(t^{2.0539221767\ldots})$. A finite-fiber construction gives the range $s^2(\log s)^2=o(\log t)$. For fixed parameters, algebraic preprocessing terminates and field initialization, vertex decoding and adjacency take deterministic time polynomial in the extension degree. We also retain an elimination-free variant and formulate the filtering argument for general forbidden configurations with a distinguished edge.
\end{abstract}

\noindent\textit{Keywords.} Ramsey graphs; finite fields; incident flags; polynomial interpolation; algebraic configurations.\par
\noindent\textit{2020 Mathematics Subject Classification.} Primary 05D10; Secondary 05B25, 14G15, 68R10.

\section{Introduction}\label{sec:intro}
The Ramsey number $R(s,t)$ is the least $N$ such that every graph on $N$ vertices contains a clique of size $s$ or an independent set of size $t$. A graph $G$ with $N$ vertices, clique number $\omega(G)<s$ and independence number $\alpha(G)<t$ proves $R(s,t)>N$. We study explicit witnesses in the fixed off-diagonal regime: $s$ is fixed and $t\to\infty$. Unless another base is displayed, logarithms are to base two.

Ihringer and Mattheus \cite{IM} construct ordered trace-ratio graphs over $\F_{2^h}$ with independence number at most $d^h+1$ and clique number at most $2^{d-2}$. Their incidence-rank argument extends to any family of incident point--hyperplane pairs. We choose those flags by polynomial interpolation, so that the rank bound holds for every coefficient assignment while the coefficients remain free to prevent cliques. The two tasks impose different constraints: independence is controlled by the point features, and clique avoidance is controlled by the dimension of ordered configurations.

\subsection{The Ramsey bounds}
For $3\le d\le s$, define
\[
\kappa_{s,d}=s(d-1)-\binom d2,\qquad
\mu_{s,d}=\frac{\kappa_{s,d}}{s-1},\qquad
a=\lceil\mu_{s,d}\rceil,\quad r=d-1-a,\quad k=\lceil\log_2s\rceil.
\]
If $r=0$, put $\rho=d$. Otherwise let $\rho$ be the root in $(d-1,d)$ of $x^k-dx^{k-1}+r=0$. Set
\[
E_{\rm crit}(s)=\max_{3\le d\le s}\frac{\mu_{s,d}}{\log_2d},\qquad
E_{\rm F}(s)=\max_{3\le d\le s}\frac{\mu_{s,d}}{\log_2\rho}.
\]
The first exponent uses the ordinary rank certificate; the second uses short Frobenius relations between point coordinates.

\begin{theorem}\label{thm:main}
For every fixed integer $s\ge3$, deterministic witnessing graphs give
\[
R(s,t)\ge\Omega_s\bigl(t^{E_{\rm F}(s)}\bigr),\qquad
E_{\rm F}(s)\ge E_{\rm crit}(s)\ge\frac{s}{2\log_2s}.
\]
Both optimized exponents are $(1/2+o(1))s/\log_2s$ as $s\to\infty$. For fixed parameters, algebraic preprocessing terminates and field initialization, vertex decoding and adjacency have bit cost polynomial in the extension degree.
\end{theorem}

The construction and extraction are proved in Theorem~\ref{thm:critical}, Proposition~\ref{prop:two-stage} and Corollary~\ref{cor:frobenius-exponent}; Proposition~\ref{prop:asymptotic} gives the asymptotics. The notation $\Omega_s$ includes a positive prefactor and an availability threshold depending on $s$. These constants are independent of $t$.

For $s=16$, take $d=13$ and write $\mu=38/5$. With $\rho_{16}$ the root in $(12,13)$ of $x^4-13x^3+4=0$, the bound is
\[
R(16,t)\ge\Omega\bigl(t^{38/(5\log_2\rho_{16})}\bigr)
=\Omega\bigl(t^{2.0539221767\ldots}\bigr).
\]
The unrefined critical exponent is already $2.0538099736\ldots$. Thus the endpoint and removal of the integral rounding loss account for the main gain in this example; the short relation improves the denominator further.

\begin{table}[t]
\centering\small
\setlength{\tabcolsep}{5pt}
\begin{tabular}{rccccc}
\toprule
$s$ & $E_{\rm dir}$ & $E_{\rm mark}$ & $E_{\rm crit}$ & $E_{\rm F}$ & Prior benchmark\\
\midrule
4  & 0.630930 & 0.630930 & 1.051550 & 1.052649 & $5/3$\\
6  & 0.861353 & 1.000000 & 1.205894 & 1.212043 & $2$\\
12 & 1.505150 & 1.666667 & 1.724081 & 1.724306 & $2$\\
15 & 1.734389 & 1.952601 & 1.972525 & 1.972640 & $2$\\
16 & 1.891667 & 1.952601 & 2.053810 & 2.053922 & $2$\\
17 & 1.952601 & 2.101196 & 2.134027 & 2.134033 & $2$\\
18 & 2.047664 & 2.161905 & 2.213284 & 2.213289 & $2$\\
33 & 3.191191 & 3.328234 & 3.334734 & 3.334734 & 2.137243\\
\bottomrule
\end{tabular}
\caption{Optimized sufficient exponents, rounded to six decimals. The original strict constructions give $E_{\rm dir}$ and $E_{\rm mark}$ in \eqref{eq:old-optimized}; the endpoint constructions give $E_{\rm crit}$ and $E_{\rm F}$. The final column takes the larger of the Ihringer--Mattheus exponent and the finite-geometric benchmarks discussed below. It is a comparison of these cited constructions, not an exhaustive record claim.}\label{tab:comparison}
\end{table}

\subsection{Previous constructions and the computational scope}
Alon and Pudl\'ak \cite{AP} give explicit exponents of order $\sqrt{\log s/\log\log s}$. The uniform Ihringer--Mattheus bound \cite[Theorem 1.3]{IM} has exponent $(1-o(1))\lceil\log_2s\rceil/\log_2(\lceil\log_2s\rceil+1)$ for $3\le s\le t$. The fixed-$s$ scale here is larger as $s\to\infty$, but our assertion of polynomial local access fixes the algebraic parameters first.

Small clique sizes require a separate comparison. Kostochka, Pudl\'ak and R\"odl \cite{KPR} give exponents $8/5$ for $s=4$, $5/3$ for $s=5$, and $2$ for $s=6$; the last applies to every $s\ge6$ by monotonicity. Han, Ihringer and Van Maldeghem \cite{HIV} improve the explicit $s=4$ exponent to $5/3$ using pseudorandom graphs in the split Cayley hexagon. The flag-family bounds are weaker than these benchmarks for small $s$. Within the parameter choices analyzed here, the refined endpoint first exceeds $2$ at $s=16$, the original marked bound at $s=17$, and the original direct bound at $s=18$. Table~\ref{tab:comparison} also includes the Ihringer--Mattheus benchmark at $s=33$.

These comparisons concern explicit witnessing constructions. For unrestricted witnesses, Brada\v{c} \cite{Bradac} proves $R(s,t)\ge\Omega_s(t^{s-1}/(\log t)^{2s-4})$ for fixed $s\ge3$. The smaller exponent here is accompanied by deterministic local access.

We represent $\F_{2^h}$ by binary polynomials modulo an irreducible polynomial of degree $h$. Shoup's algorithm \cite{Shoup} constructs one deterministically in time polynomial in $h$ in fixed characteristic, and every field operation has polynomial bit cost. Preprocessing computes fixed prime-field polynomials before $h$ varies. Its dependence on $s,d$ is not claimed to be efficient. Coefficient selection, decoding an indexed vertex, and testing adjacency then have cost polynomial in $h$. The resulting graphs have exponentially many vertices and are accessed locally.

There is also a quantified existence range. Corollary~\ref{cor:growing-range} gives $R(s,t)\ge t^{(1/2-o(1))s/\log_2s}$ when $s\to\infty$ and $s^2(\log s)^2=o(\log t)$. A sufficient threshold has size $2^{O(s^2(\log s)^2)}$. This bounds availability, not the preprocessing cost as $s$ varies. The original elimination-free construction, with its larger threshold and strict dimension inequality, remains in Appendix~\ref{sec:original-statements}.

\subsection{The configuration method}
In a chart of dimension $\nu$, suppose an ordered forbidden family of $v$ entries has codimension $\kappa$. Copy $m$ coordinates as labels and interpolate the remaining coordinates independently. At every distinct label tuple, universal evaluation rank gives the coefficient--configuration incidence bound $L+mv-\kappa$, where $L$ is the number of coefficients. A strict inequality $m(v-1)<\kappa$ makes the projection remembering one label proper, so deletion of a polynomial zero set suffices. At equality it may be dominant. Remembering two labels instead gives an image of dimension at most $L+m$ that cannot contain its $L+m$-dimensional diagonal: the actual tuples have distinct labels. A polynomial separates the image from that diagonal. Deleting an edge where this polynomial vanishes excludes every forbidden tuple; its nonzero diagonal lets us multiply the old triangular rank certificate by the filter.

Restriction of scalars multiplies the configuration codimension by $b$, allowing $m/b=\kappa/(v-1)$ exactly. The following principle isolates the inputs used by the method; its proof is in Section~\ref{sec:general}.

\begin{theorem}[Distinguished-edge configuration principle]\label{thm:configuration}
Let a prime-field polynomial chart have dimension $\nu$ and an ordered graph with an assignment-uniform triangular rank certificate. Let $\Ccal_i\subseteq(\A^\nu)^{v_i}$, with $v_i\ge2$, for finitely many $i$, be constructible forbidden configuration families of codimension at least $\kappa_i$. Require distinct copied labels and choose a distinguished edge of each family. After restriction of scalars of degree $b$, use $m$ copied base coordinates and independent interpolating outputs. If
\[
m\le b\nu,\qquad m(v_i-1)\le b\kappa_i\quad\text{for every }i,
\]
then finite deterministic preprocessing, specialization and polynomial edge filtering produce a positive-density indexed vertex set avoiding all these patterns. The independence bound changes by a fixed multiplier. For fixed model parameters, decoding and adjacency have cost polynomial in the base-field extension degree.
\end{theorem}

For flags, the exact prefix-rank calculation includes degenerate configurations and proves the codimension is sharp. Compressed interpolation works at every distinct tuple, including exceptional labels. Degree control, prime-field descent and specialization make the resulting polynomial prescription deterministic. The dense decoder supplies an index set rather than only a membership test. These details are needed to turn the dimension argument into locally accessible graphs.

The relation to evasive sets \cite{DKL,LNZ} lies in controlling intersections with algebraic configurations. Those works bound how a set meets varieties; here the forbidden object is an ordered tuple in a polynomially parameterized graph, and the rank estimate is independent of its coefficients. Section~\ref{sec:general} also gives the fixed-characteristic version. Appendices contain the decoder, the original strict variants, the quantified finite-fiber argument, and the elementary geometric conventions.
% END SOURCE_BLOCK revised_intro

% BEGIN SOURCE_BLOCK chart_rank
\section{The graph and its rank certificate}\label{sec:rank}
Let $\Bbbk$ be a field and put $n=d-1$. A projective point in $\PP^{d-1}(\Bbbk)$ is a one-dimensional subspace of $\Bbbk^d$. A nonzero column vector represents it, with scalar multiples representing the same point. A hyperplane is the kernel of a nonzero linear functional, represented by a row vector up to nonzero scale. A flag is an incident pair: its point lies in its hyperplane, or equivalently the row applied to the column is zero.

We use points whose first coordinate is nonzero and hyperplanes whose last coefficient is nonzero. Scaling makes these coordinates equal to $1$. Write
\begin{equation}\label{eq:chart}
P(p)=(1,p_1,\ldots,p_n)^T,\qquad
H(p,\beta)=\left(-p_n-\sum_{i=1}^{n-1}\beta_i p_i,\beta_1,\ldots,\beta_{n-1},1\right).
\end{equation}
The free parameters are $p\in\Bbbk^n$ and $\beta\in\Bbbk^{n-1}$. There are $2d-3$ of them, so this chart of the flag space is identified with $\A^{2d-3}$. Its first hyperplane coefficient is forced by incidence:
\[
H(p,\beta)P(p)=-p_n-\sum_{i=1}^{n-1}\beta_i p_i+\sum_{i=1}^{n-1}\beta_i p_i+p_n=0.
\]
Both representing vectors are nonzero because of their normalized coordinates. Self-incidence is built into the parameterization.

For a label $z=(z_1,\ldots,z_m)$ in the original construction, copy $p_i(z)=z_i$ for $1\le i\le m$. The remaining $r=n-m$ point coordinates and all $d-2=n-1$ slope coordinates are polynomial functions of $z$. There are $c=r+d-2$ outputs. Initially their coefficients are independent variables; their values will be fixed later. The point attached to a label recovers that label from its copied coordinates, so different labels give different projective points. Hyperplanes may repeat.

On a finite label set choose a total order $\prec$. Distinct labels $z\prec w$ are adjacent precisely when
\begin{equation}\label{eq:adjacency}
z\sim w\quad\Longleftrightarrow\quad K(z,w)=H(w)P(z)=0,
\end{equation}
where
\begin{equation}\label{eq:kernel}
K(z,w)=p_n(z)-p_n(w)+\sum_{i=1}^{n-1}\beta_i(w)\bigl(p_i(z)-p_i(w)\bigr).
\end{equation}
This tests whether the earlier point belongs to the later hyperplane. The graph is undirected because the order chooses the test for each unordered pair; $K$ itself need not be symmetric. Equal labels are never joined. In characteristic two the minus signs equal plus signs.

\begin{lemma}\label{lem:rank}
For flags in the chart \eqref{eq:chart} over $\F_{2^H}$, the graph \eqref{eq:adjacency} has independence number at most $d^H$, for every assignment and every total order. More precisely, if
\begin{equation}\label{eq:feature-space}
\Wcal_H(P)=\spanop\left\{
\prod_{\ell=0}^{H-1}P_{\sigma(\ell)}^{2^\ell}:\
\sigma:\{0,\ldots,H-1\}\longrightarrow\{0,\ldots,d-1\}
\right\}
\end{equation}
is a space of functions on the labels, then $\alpha\le\dim\Wcal_H(P)$. Restricting the label set preserves these bounds.
\end{lemma}

\begin{proof}
For this proof, $P_i(z)$ and $H_i(w)$ denote coordinates indexed by $0\le i\le d-1$. Put $Q=2^H$ and form
\[
B_{z,w}=1-K(z,w)^{Q-1}.
\]
Every nonzero element of $\F_Q$ has $(Q-1)$st power $1$, whereas $0^{Q-1}=0$. Thus $B$ is the incidence indicator, and self-incidence gives $B_{z,z}=1$.

If $z_1\prec\cdots\prec z_u$ is independent, then $K(z_i,z_j)\ne0$ for $i<j$, so $B_{z_i,z_j}=0$. Its principal submatrix is lower triangular with diagonal entries $1$. Hence
\begin{equation}\label{eq:triangular}
u\le\rank_{\F_Q}B.
\end{equation}
In characteristic two, $(x+y)^{2^\ell}=x^{2^\ell}+y^{2^\ell}$. Since $Q-1=1+2+\cdots+2^{H-1}$,
\begin{equation}\label{eq:frobenius-factor}
K(z,w)^{Q-1}=
\prod_{\ell=0}^{H-1}\left(\sum_{i=0}^{d-1}P_i(z)^{2^\ell}H_i(w)^{2^\ell}\right)
=\sum_\sigma U_\sigma(z)V_\sigma(w),
\end{equation}
where
\[
U_\sigma(z)=\prod_{\ell=0}^{H-1}P_{\sigma(\ell)}(z)^{2^\ell},\qquad
V_\sigma(w)=\prod_{\ell=0}^{H-1}H_{\sigma(\ell)}(w)^{2^\ell}.
\]
Each summand is an outer product. Let $\sigma_0$ be constantly zero. The normalized coordinate $P_0=1$ gives $U_{\sigma_0}=1$, so the constant term is absorbed into the same feature:
\[
B_{z,w}=U_{\sigma_0}(z)\bigl(1-V_{\sigma_0}(w)\bigr)
-\sum_{\sigma\ne\sigma_0}U_\sigma(z)V_\sigma(w).
\]
Every column of $B$ therefore lies in $\Wcal_H(P)$, a space spanned by at most $d^H$ functions. Together with \eqref{eq:triangular} this proves the bounds. Restriction only restricts the evaluations in this factorization.
\end{proof}

The degrees of the point and hyperplane coordinates do not enter the proof. We can give the polynomial family enough freedom to exclude cliques without weakening the independence estimate. This is the incidence-rank method of \cite{IM}, with the affine constant feature accounted for explicitly.
% END SOURCE_BLOCK chart_rank

% BEGIN SOURCE_BLOCK geometry_setup
\section{The geometry of every ordered clique}\label{sec:geometry}
Let $\Ccal_{s,d}$ consist of ordered incident point--hyperplane flags satisfying
\begin{equation}\label{eq:clique-config}
P_i\subseteq H_j\qquad(1\le i<j\le s).
\end{equation}
Positions give the configuration order; there is no order on the algebraically closed field. Repeated points and flags are permitted. Every clique supplies such a tuple, so a dimension bound for this larger family applies to the cliques to be excluded.
% END SOURCE_BLOCK geometry_setup

% BEGIN SOURCE_BLOCK exact_geometry
\subsection{Exact prefix-rank strata}\label{sec:exact}
All dimensions in this section are over $\overline{\F}_2$. Put
\begin{equation}\label{eq:kappa}
u_* =\min(s,d-1),\qquad
\kappa_{s,d}=s(d-1)-\frac{u_*(2d-u_*-1)}2
=\begin{cases}
s(d-1)-\binom d2,&d\le s+1,\\
\binom s2,&d\ge s+1.
\end{cases}
\end{equation}
The two formulas agree at $d=s+1$.

\begin{theorem}\label{thm:exact-dimension}
The space of ordered incident flags satisfying \eqref{eq:clique-config} has dimension
\[
\dim\Ccal_{s,d}=s(d-2)+\frac{u_*(2d-u_*-1)}2.
\]
The same dimension is attained in the affine chart \eqref{eq:chart}. Thus its codimension in $s$ copies of that chart is $\kappa_{s,d}$.
\end{theorem}

\begin{proof}
For each position put $R_j=\spanop(P_1,\ldots,P_j)$, with $R_0=0$. Ordered incidence together with self-incidence gives $R_j\subseteq H_j$, so every prefix span is proper. Stratify by its rank path $\rho_j=\dim R_j$, ending at rank $u$. The first rank is $1$; every subsequent step preserves the rank or increases it by one.

Fix the choices through position $j-1$. If the rank stays at $v=\rho_j$, the new point is in $\PP(R_{j-1})$, of dimension $v-1$. After choosing it, $H_j$ can be any hyperplane containing $R_j$. Such hyperplanes form the projective dual of $\Bbbk^d/R_j$, of dimension $d-v-1$. The step therefore contributes $d-2$.

If the rank increases to $v$, the new point is in $\PP^{d-1}\setminus\PP(R_{j-1})$, an open set of dimension $d-1$. The hyperplane still contributes $d-v-1$, for a total of $(d-2)+(d-v)$. This includes the first step, from rank $0$ to rank $1$. There is exactly one increase to each rank $v=1,\ldots,u$, giving
\begin{equation}\label{eq:rank-stratum}
s(d-2)+\sum_{v=1}^u(d-v)=s(d-2)+\frac{u(2d-u-1)}2.
\end{equation}
These are dimension equalities, not just parameter upper bounds. On a finite cover where the prefix bases are chosen by nonzero minors, the successive point choices are projective-space bundles or open subsets of them, and the hyperplane choices are projective-space bundles. Earlier hyperplanes do not constrain later points. The dimensions therefore add. The expression increases with $u$ up to $u_*$. This final rank is realized by taking the first $u_*$ points independent and later points in their span.

To see that a maximal stratum meets the affine chart, choose its final span not to contain the last coordinate vector $e_{d-1}$ and to contain a vector with nonzero first coordinate. Within it choose points with the required ranks and nonzero first coordinate. Every prefix then admits an annihilating functional nonzero on $e_{d-1}$: otherwise $e_{d-1}$ would lie in that prefix. Normalize this functional to have last coefficient $1$. These conditions define a nonempty open part of the stratum, so the chart has the same dimension. Subtracting \eqref{eq:rank-stratum} at $u=u_*$ from $s(2d-3)$ gives \eqref{eq:kappa}.
\end{proof}

For $d\le s+1$, the complete-flag bound is therefore exact. In the original copied-coordinate model with $m\ge1$, the maximal stratum also has a nonempty distinct-label open subset: its point span can be chosen so that projection onto the first $m$ affine coordinates is nonconstant, after which finitely many pairwise coincidences are avoided generically. Counting the same unrestricted configuration space more carefully cannot lower its dimension in the range relevant to the main exponent.
% END SOURCE_BLOCK exact_geometry

% BEGIN SOURCE_BLOCK interpolation
\section{Universal interpolation and the incidence dimension}\label{sec:interpolation}
We need to prescribe each polynomial output at any $s$ distinct labels. Restricting to generic labels would leave exceptional cliques untreated. For $z\in\Bbbk^m$ and a formal indeterminate $T$, put
\begin{equation}\label{eq:compressed-features}
\ell_T(z)=z_1+z_2T+\cdots+z_mT^{m-1},\qquad
\phi_{j,e}(z)=[T^e]\ell_T(z)^j
\end{equation}
for $0\le j<s$ and $0\le e\le(m-1)j$. Each feature has label degree at most $j$. Their number of slots is
\begin{equation}\label{eq:M}
M(m,s)=\sum_{j=0}^{s-1}\bigl((m-1)j+1\bigr)=s+(m-1)\binom s2.
\end{equation}
Slots need not be independent. For instance, in characteristic two $(z_1+z_2T)^2=z_1^2+z_2^2T^2$, so $\phi_{2,1}=0$ when $m=2$. We retain every slot; what matters is evaluation rank.

\begin{lemma}[Universal interpolation]\label{lem:interpolation}
For any field and any $s$ pairwise distinct labels $z_1,\ldots,z_s$, the $s\times M(m,s)$ matrix
\[
E(z_1,\ldots,z_s)=\bigl(\phi_{j,e}(z_i)\bigr)_{i,(j,e)}
\]
has rank $s$. Linear combinations of the features can take arbitrary prescribed values at those labels.
\end{lemma}

\begin{proof}
First work over an algebraic closure. For $i\ne i'$, the polynomial $\ell_T(z_i)-\ell_T(z_{i'})$ is nonzero, since its coefficients are coordinate differences. Choose $\tau$ outside the finitely many roots of all such polynomials. Then $a_i=\ell_\tau(z_i)$ are pairwise distinct. For each $j<s$,
\[
\ell_\tau(z)^j=\sum_{e=0}^{(m-1)j}\tau^e\phi_{j,e}(z).
\]
The evaluation columns of $1,\ell_\tau,\ldots,\ell_\tau^{s-1}$ belong to the column span of $E$. They form a Vandermonde matrix with determinant $\prod_{i<i'}(a_{i'}-a_i)\ne0$. Hence $E$ has rank $s$. Matrix rank is unchanged by extending scalars, because nonvanishing of a minor descends. The auxiliary $\tau$ need not belong to the field of the labels.
\end{proof}

Enumerate slots by increasing $j$, then increasing $e$. Give each output its own independent block. In the original flag model the output order is $p_{m+1},\ldots,p_n,\beta_1,\ldots,\beta_{n-1}$. This fixes an order on the $L=cM(m,s)$ coefficients $A=(A_0,\ldots,A_{L-1})$. Distinct coefficient tuples remain distinct in parameter space even if they give the same polynomial functions.

\begin{lemma}\label{lem:universal-dimension}
Let $\Ccal\subseteq(\A^\nu)^s$ be constructible, with $\dim\Ccal\le\nu s-\kappa$. Copy $m$ chart coordinates as labels, require pairwise distinct labels, and parameterize the other $c=\nu-m$ coordinates by independent blocks of the features \eqref{eq:compressed-features}. The coefficient--configuration incidence set then satisfies
\begin{equation}\label{eq:universal-dimension}
\dim\Ical\le L+ms-\kappa.
\end{equation}
\end{lemma}

\begin{proof}
Fix a full configuration, including its labels and output values. For one output the evaluation conditions have matrix $E$ and rank $s$, by Lemma~\ref{lem:interpolation}. They are onto and leave an affine solution space of dimension $M-s$. Disjoint coefficient blocks leave dimension $L-cs$ in total.

Cover the configuration family by the finitely many open sets on which a specified $s\times s$ evaluation minor is nonzero. On one such chart, leave $M-s$ coefficients free in each block and solve for the other $s$. Matrix inversion has only powers of the chosen nonzero determinant in its denominators, so the solved coefficients are regular on the chart. The incidence is a product there with $\A^{L-cs}$. Dimensions add, and the finite cover gives
\[
\dim\Ical\le\dim\Ccal+L-cs\le L+ms-\kappa.
\]
\end{proof}

In the original model, let $\Ccal_{s,d}^{\circ}$ be the ordered configurations in the chart with pairwise distinct first $m$ point-coordinate labels. The incidence $\Ical$ records coefficients that assign exactly each flag at its own label. Lemma~\ref{lem:universal-dimension} gives
\begin{equation}\label{eq:original-dimension}
\dim\Ical\le L+\binom d2-rs.
\end{equation}
The coefficient-image closure is proper if $rs>\binom d2$. The image closure remembering also one label is proper if $\binom d2-rs<m$, equivalently $r(s-1)>\binom{d-1}2$. Distinctness is essential: different flags with the same copied label need not impose independent evaluation conditions.
% END SOURCE_BLOCK interpolation

% BEGIN SOURCE_BLOCK degree_tools
\section{Degree, descent and effective specialization}\label{sec:algebra}
Properness gives a nonzero equation for an exceptional set. A deterministic coefficient prescription also needs a degree bound and an equation over the prime field. We use cumulative degree to retain components of every dimension. For a reduced affine algebraic set $X$, set
\[
\cdeg X=\sum_{Z\text{ an irreducible component of }X}\deg Z,
\qquad \cdeg\varnothing=0.
\]
The degree of an affine irreducible set is the degree of its projective closure. A generic complementary affine linear section consists of that many points. In particular, a finite reduced set has cumulative degree equal to its number of geometric points.

\begin{lemma}\label{lem:degree}
The following statements hold over an algebraically closed field.
\begin{enumerate}
\item If $f_1,\ldots,f_u$ have positive degree, then
\begin{equation}\label{eq:bezout}
\cdeg\bigl(X\cap V(f_1,\ldots,f_u)\bigr)\le\cdeg X\prod_{i=1}^u\deg f_i.
\end{equation}
\item If $\pi$ is linear, then $\cdeg\overline{\pi(X)}\le\cdeg X$.
\item If $Y\subsetneq\A^N$ is proper, there is a nonzero polynomial of degree at most $\cdeg Y$ vanishing on $Y$. For $Y=\varnothing$, use $1$.
\item Over $\overline{\F}_2$, if $Y$ is invariant under coordinatewise Frobenius and its inverse, a nonzero vanishing polynomial with a given degree bound can be chosen over $\F_2$ without increasing that bound.
\end{enumerate}
\end{lemma}

\begin{proof}
The first assertion is cumulative B\'ezout \cite[Theorem 1.1]{LR}. Intersect one irreducible component with one hypersurface at a time. If the component is contained in the hypersurface, its degree is unchanged and at most its old degree times the hypersurface degree. Otherwise projective B\'ezout bounds the sum of intersection-component degrees by the product. Summing, possibly overcounting coincident pieces, and iterating gives \eqref{eq:bezout}. Zero equations may be omitted; a nonzero constant equation gives the empty set. No equidimensionality or complete-intersection assumption is used.

For the projection assertion, first let $X$ be irreducible of dimension $e$ and let $Y=\overline{\pi(X)}$ have dimension $f$. By constructibility and the fiber-dimension theorem \cite[Chapter I]{Hartshorne}, the image contains a nonempty open subset of $Y$ over which fibers have dimension $e-f$. A generic complementary affine section of $Y$ consists of $\deg Y$ distinct points in that open subset and avoids infinity. Pull back its $f$ linear equations and add $e-f$ generic source hyperplanes. On each of the finitely many selected fibers, these meet in a finite nonempty set, avoiding its boundary at infinity. The resulting finite intersection in the source has at least $\deg Y$ points and is cut out by $e$ linear equations. Equation~\eqref{eq:bezout} bounds its size by $\deg X$. For reducible $X$, sum over components; image components are among the closures of these component images.

For the third assertion, take an irreducible component $Z$ of dimension $f<N$. A generic surjective linear map $\lambda:\A^N\to\A^{f+1}$ has image closure of dimension $f$ on $Z$. This follows from generic linear Noether normalization, adding one independent ambient linear form. The image closure is an irreducible hypersurface, given by an irreducible polynomial $g_Z$. The projection assertion gives $\deg g_Z\le\deg Z$. Its pullback is nonzero, since $\lambda$ is surjective, and vanishes on $Z$. Multiply the pullbacks over components. The product is nonzero and has degree at most $\cdeg Y$.

For descent, list all monomials up to the required degree. Their coefficient vectors for polynomials vanishing on $Y$ form a nonzero $\overline{\F}_2$-linear space $V$. If $f=\sum c_\nu X^\nu$ vanishes on $Y$, its coefficientwise square $f^{(2)}=\sum c_\nu^2X^\nu$ does too: for $y\in Y$, its coordinatewise square root $x$ belongs to $Y$ and $f^{(2)}(y)=f(x)^2=0$. Frobenius and its inverse therefore preserve $V$.

Reduce a basis matrix of $V$ to reduced row-echelon form. Squaring its entries gives another reduced matrix with the same row space and pivot pattern. Uniqueness makes the two matrices identical. Every entry satisfies $c^2=c$, hence lies in $\F_2$. Any nonzero row gives the required prime-field polynomial. The degree bound is unchanged.
\end{proof}

A prime-field zero set and a coordinate projection commute with Frobenius. Frobenius is a homeomorphism, so projected-image closures have the invariance needed for descent. Individual geometric components may be permuted; the full exceptional set is what descends.
% END SOURCE_BLOCK degree_tools

% BEGIN SOURCE_BLOCK adaptive_degree
\subsection{A smaller model after substitution}\label{sec:degree-model}
\begin{proposition}\label{prop:substituted-degree}
In a fixed degree-$b$ restriction-of-scalars model, suppose all substituted point and slope coordinates have total degree at most $e$ in coefficients and labels. The closure $X$ of the clique incidence, restricted by distinctness open conditions, satisfies
\begin{equation}\label{eq:substituted-degree}
\cdeg X\le\mathcal D:=(2e)^{b\binom s2}.
\end{equation}
For a single-stage compressed family, one may take $e=s$. For the two-stage family of Section~\ref{sec:two-stage}, one may take $e=1+s(s-1)$.
\end{proposition}

\begin{proof}
The expression \eqref{eq:kernel} has degree at most $2e$ after substitution. Multiplication in the fixed degree-$b$ algebra is bilinear in its base coordinates with fixed structure constants. Each clique pair therefore supplies $b$ scalar equations of degree at most $2e$. Cumulative B\'ezout bounds the full zero set $V$ by $(2e)^{b\binom s2}$.

The distinctness requirements form an ambient open set $O$. The closure of $V\cap O$ is precisely the union of the irreducible components of $V$ meeting $O$: the nonempty open intersection with any such component is dense in it. Passing to that union cannot increase cumulative degree.
For the single-stage family, a feature has label degree at most $s-1$ and its coefficient contributes at most one, so every substituted chart coordinate has degree at most $s$. In the two-stage family, each base coordinate of $X$ has degree at most $s$ in the first-stage coefficients and labels. A second-stage feature has degree at most $s-1$ in $X$; its coefficient adds one. The resulting total degree is at most $1+s(s-1)$. Bilinear algebra multiplication does not alter these bounds on the composed coordinate polynomials.
\end{proof}

This bound requires no inverse-witness variables. Such variables remain useful for computation. To express $z_i\ne z_j$, introduce $u_{ij,1},\ldots,u_{ij,m}$ and impose
\begin{equation}\label{eq:inverse-witness}
\sum_{\ell=1}^m u_{ij,\ell}(z_{i,\ell}-z_{j,\ell})=1.
\end{equation}
The equation is solvable exactly when one coordinate difference is nonzero. For a candidate polynomial $Q$ on a projected incidence image, append $vQ-1=0$. The candidate vanishes on that image, and hence its closure, exactly when the enlarged system has no algebraically closed point. By the Nullstellensatz this is equivalent to its ideal containing $1$. A Gr\"obner-basis computation over $\F_2$ decides this and terminates. Extending the coefficient field does not change whether an ideal is the unit ideal. These witness variables are used for ideal membership, never to establish the incidence dimension.

\subsection{Specializing a computed polynomial}\label{sec:adaptive}
\begin{lemma}\label{lem:adaptive}
Let $R\in\F_q[A_0,\ldots,A_{L-1}]$ be nonzero of total degree at most $\delta$, with $q>\delta$. A tuple $a\in\F_q^L$ satisfying $R(a)\ne0$ can be found deterministically by testing at most $\delta+1$ fixed field elements at each coordinate and maintaining a nonzero partially specialized polynomial.
\end{lemma}

\begin{proof}
Suppose the first coordinates have been fixed and the remaining polynomial is nonzero. Regard it as a polynomial in the variables after $A_i$, with coefficients that are polynomials in $A_i$. One coefficient $g(A_i)$ is nonzero and has degree at most $\delta$. A value that makes the entire remaining polynomial zero must be a root of $g$. At most $\delta$ values do this, so the first $\delta+1$ distinct field elements contain a choice preserving nonzeroness. Continue. The final nonzero polynomial in no variables is a nonzero constant.
\end{proof}

If $Q(A,Z)=\sum_\gamma R_\gamma(A)Z^\gamma$ is nonzero, choose its first nonzero coefficient polynomial and apply the lemma. The resulting $Q(a,Z)$ is nonzero. For fixed degrees and dimensions, the symbolic coefficient comparisons and field operations have bit cost polynomial in $h$ when $q=2^h$. Once $Q$ is computed, the sufficient field-size condition is $q>\delta$, rather than $h\ge(\delta+1)^L$. To retain at least half the labels by a degree-$\delta$ deletion, $q>2\delta$ suffices.
% END SOURCE_BLOCK adaptive_degree

% BEGIN SOURCE_BLOCK critical_filter
\section{Polynomial edge deletion at the critical dimension}\label{sec:filter}
At the strict marked inequality, the coefficient--label projection is proper and one can delete a vertex zero set. At equality that projection may be dominant. Projecting instead to the first two labels produces a correspondence that can be separated from the diagonal. Its vanishing polynomial removes one edge of every forbidden clique and still gives a rank certificate on the retained vertices.

\begin{lemma}\label{lem:filter-rank}
Let an ordered graph have a matrix $B$ over a field with nonzero diagonal and $B_{z,w}=0$ on every ordered nonedge $z\prec w$. Let
\[
F(Z,W)=\sum_{i=1}^R f_i(Z)g_i(W).
\]
Retain vertices with $F(z,z)\ne0$, and retain an old edge $z\prec w$ only if $F(z,w)\ne0$. Then the modified graph satisfies $\alpha\le R\rank B$. A polynomial with $\deg_WF\le\delta$ admits $R\le\binom{m+\delta}{m}$.
\end{lemma}

\begin{proof}
The matrix $M=B\circ(F(z,w))_{z,w}$, where $\circ$ denotes entrywise multiplication, has nonzero diagonal. On an independent set of the modified graph, an entry above the diagonal is zero: either the old pair was a nonedge, making $B_{z,w}=0$, or its edge was deleted, making $F(z,w)=0$. Thus the corresponding submatrix is triangular and nonsingular.

If $B=\sum_{j=1}^{R_B}u_jv_j^T$ with $R_B=\rank B$, then
\[
M=\sum_{j=1}^{R_B}\sum_{i=1}^R(u_j\circ f_i)(v_j\circ g_i)^T.
\]
It has rank at most $R_BR$. Expansion in $W$-monomials gives the binomial bound.
\end{proof}

\begin{theorem}[Diagonal separation]\label{thm:diagonal-separation}
Let $\Ical\subseteq\A^L\times(\A^m)^s$ be a prime-field constructible family of forbidden ordered tuples with pairwise distinct labels. Suppose
\begin{equation}\label{eq:critical-dimension}
\dim\Ical\le L+m.
\end{equation}
There is a polynomial $F(A,Z,W)\in\F_2[A,Z,W]$ that vanishes on the first two labels of every tuple in $\Ical$ and has
\begin{equation}\label{eq:nonzero-diagonal}
D_F(A,Z):=F(A,Z,Z)\not\equiv0.
\end{equation}
It can be computed by terminating deterministic preprocessing from a fixed polynomial incidence model. If positions $1,2$ are an edge of every forbidden pattern and the old graph has an assignment-uniform triangular rank certificate, coefficient specialization, vertex deletion and filtering by $F\ne0$ exclude every such pattern at a fixed multiplicative independence cost. For fixed model parameters, decoding and adjacency are polynomial in $h$ over $\F_{2^h}$.
\end{theorem}

\begin{proof}
Write $S$ for the actual image of $\Ical$ under $(A,z_1,\ldots,z_s)\mapsto(A,z_1,z_2)$, and let $Y=\overline S$. Let
\[
D=\{(A,Z,W):Z=W\}\subseteq\A^{L+2m}.
\]
The diagonal is irreducible of dimension $L+m$, whereas $\dim Y\le L+m$. Although $S$ is disjoint from $D$, its closure may meet $D$. We must show that it cannot contain the whole diagonal.

First, $D$ cannot be an irreducible component of $Y$. If it were, write the other components as $Y_1,\ldots,Y_u$. Since none contains $D$, the set $D\setminus\bigcup_iY_i$ is nonempty. It equals $Y\setminus\bigcup_iY_i$, so is open in $Y$. It is disjoint from $S$, contradicting the density of $S$ in $Y$. Next, if $D\subseteq Y$, irreducibility places it in a single component $Y_i$. Since $\dim Y_i\le\dim D$, the inclusion forces $D=Y_i$, which has just been excluded. Hence $D\nsubseteq Y$.

Choose a polynomial $f\in I(Y)$ whose restriction to $D$ is nonzero. The space of polynomials in $I(Y)$ of degree at most $\deg f$ is invariant under coefficient Frobenius and its inverse, because the model and the projection are defined over $\F_2$. Lemma~\ref{lem:degree} gives a basis over $\F_2$ by row-echelon descent. If every basis member vanished on $D$, so would $f$. At least one therefore supplies the required $F\in\F_2[A,Z,W]$.

To compute it, enumerate total degrees and then binary coefficient vectors in a fixed order. Reject a candidate whose diagonal substitution is the zero polynomial. Test vanishing on $S$ by adding $vF(A,z_1,z_2)-1=0$ to an inverse-witness model for $\Ical$ and checking whether the ideal contains $1$, as in Section~\ref{sec:degree-model}. The preceding existence proof ensures that the enumeration terminates. This procedure gives a finite degree depending only on the fixed model. It does not give a prescribed bound on that degree; the quantified construction in Appendix~\ref{sec:fibers} uses a different argument.

Expand $D_F(A,Z)$ in $Z$ and choose a nonzero coefficient polynomial. Lemma~\ref{lem:adaptive} specializes $A=a$ while keeping that coefficient nonzero. Thus $D_F(a,Z)$ is a nonzero polynomial. If $\delta\ge1$ bounds the computed filter's total degree, Lemma~\ref{lem:decoder} retains exactly $q^{m-1}(q-\delta)$ labels with $D_F(a,z)\ne0$, for sufficiently large $q$.

On these labels retain an old edge $z\prec w$ only when $F(a,z,w)\ne0$. Any remaining forbidden tuple would belong to $\Ical$ and hence have $F(a,z_1,z_2)=0$. Its distinguished edge would have been deleted. The retained vertices have nonzero diagonal filter values, so Lemma~\ref{lem:filter-rank} applies and gives the fixed independence multiplier. Computing the fixed polynomials, decoding labels and evaluating the old incidence test have bit cost polynomial in $h$ once the model is fixed.
\end{proof}

The geometric condition actually used is $D\nsubseteq Y$. It suffices even without \eqref{eq:critical-dimension}. Conversely, properness of $Y$ alone is insufficient: its ideal might vanish identically on $D$. A weaker inequality that merely makes the pair projection proper does not justify the critical endpoint or an endpoint beyond it.
% END SOURCE_BLOCK critical_filter

% BEGIN SOURCE_BLOCK scalars
\section{Restriction of scalars and the Ramsey endpoint}\label{sec:scalars}
Fix $b\ge1$ and an irreducible $g_b(U)\in\F_2[U]$ of degree $b$. Put $E_0=\F_2[U]/(g_b)$ and, for $q=2^h$, use
\begin{equation}\label{eq:fixed-algebra}
E_q=E_0\otimes_{\F_2}\F_q=\F_q[U]/(g_b(U)).
\end{equation}
Multiplication in the basis $1,U,\ldots,U^{b-1}$ is bilinear, with fixed binary structure constants.

\begin{lemma}\label{lem:scalars}
After base change to $\overline{\F}_2$, the algebra \eqref{eq:fixed-algebra} is the product of $b$ copies of $\overline{\F}_2$. Over $\F_{2^h}$ it is the field $\F_{2^{bh}}$ when $\gcd(b,h)=1$. The restriction of scalars of the chart has dimension $b(2d-3)$, and its ordered clique configurations have codimension $b\kappa_{s,d}$.
\end{lemma}

\begin{proof}
The polynomial $g_b$ has $b$ distinct roots, since finite fields are perfect. The Chinese remainder theorem gives the product over those roots after algebraic closure. Polynomial equations in the algebra become one independent copy in each factor, giving a product of charts and a product of configuration varieties. Dimensions and codimensions add.

A root of $g_b$ has binary Frobenius orbit of length $b$. Under $x\mapsto x^{2^h}$ its orbit has length $b/\gcd(b,h)$. Thus when $\gcd(b,h)=1$, $g_b$ remains irreducible over $\F_{2^h}$ and the quotient is the stated field.
\end{proof}

Copy $m$ base-field label coordinates into $m$ of the $b(d-1)$ base coordinates of $p$. Parameterize the remaining $c=b(2d-3)-m$ chart coordinates independently by compressed polynomials in the labels. Here $L=cM(m,s)$ counts all independent coefficient slots in the base-field description. The copied coordinates make the point map injective over the actual field $E_q$. Lemma~\ref{lem:universal-dimension} gives
\begin{equation}\label{eq:scalar-incidence}
\dim\Ical\le L+ms-b\kappa_{s,d}.
\end{equation}
Repeated points in the split geometric factors are permitted in the configuration count, so this step does not assume that distinct base labels have distinct projections to every factor.

\begin{theorem}[Critical flag construction]\label{thm:critical}
Fix $s,d\ge3$ and $b,m\ge1$ with
\begin{equation}\label{eq:endpoint-condition}
m\le b(d-1),\qquad m(s-1)\le b\kappa_{s,d}.
\end{equation}
After finite deterministic preprocessing, there are graphs for every sufficiently large $h$ with $\gcd(b,h)=1$ such that
\begin{equation}\label{eq:critical-graph}
|V|=q^{m-1}(q-\delta),\qquad\omega<s,\qquad\alpha\le C d^{bh},\qquad q=2^h.
\end{equation}
The constants $C,\delta$ depend only on the fixed parameters. Initialization, decoding and adjacency have deterministic bit cost polynomial in $h$.
\end{theorem}

\begin{proof}
The scalar family has $\dim\Ical\le L+ms-b\kappa_{s,d}\le L+m$. Describe its distinct-label incidence by the scalar equations and witnesses. Theorem~\ref{thm:diagonal-separation} computes a filter $F$ of some fixed total degree $\delta\ge1$ and specializes its coefficients so that $F(a,Z,Z)$ remains nonzero. Lemma~\ref{lem:decoder} gives exactly $q^{m-1}(q-\delta)$ vertices in that diagonal complement.

Use lexicographic order on their binary base-coordinate encodings. Join $z\prec w$ precisely when $H(w)P(z)=0$ in $E_q$ and $F(a,z,w)\ne0$ in $\F_q$. Every ordered $s$-clique would lie in the universal incidence and lose its first edge. Thus $\omega<s$.

The coprimality assumption makes $E_q=\F_{2^{bh}}$. Apply Lemma~\ref{lem:rank} over this field to get a certificate of rank at most $d^{bh}$. The scalar-valued filter is also an $E_q$-valued filter, and Lemma~\ref{lem:filter-rank} bounds its multiplier by $\binom{m+\delta}{m}$. This gives \eqref{eq:critical-graph} with a fixed $C$. The binary algebra structure is fixed, so arithmetic in $E_q$ reduces to a fixed number of base-field operations. Field initialization, specialization, decoding and the displayed adjacency rule therefore have polynomial bit cost in $h$.
\end{proof}

\begin{lemma}\label{lem:extraction}
Suppose that, for every sufficiently large $h$ in a fixed arithmetic progression of step $g$, there are $K_s$-free graphs with
\[
|V_h|\ge c_0 2^{mh},\qquad\alpha(G_h)\le C\theta^{bh},\qquad\theta>1,
\]
where all parameters and constants are fixed. Then
\begin{equation}\label{eq:bounded-gap-extraction}
R(s,t)\ge\Omega\bigl(t^{m/(b\log_2\theta)}\bigr).
\end{equation}
\end{lemma}

\begin{proof}
Choose the largest admissible $h$ with $C\theta^{bh}\le t-1$. For sufficiently large $t$ it exists in the valid range, and the gap gives
\[
h>\frac{\log_2((t-1)/C)}{b\log_2\theta}-g.
\]
Hence $|V_h|\ge c_0 2^{-mg}((t-1)/C)^{m/(b\log_2\theta)}$, while its independence number is less than $t$. Every prefactor is constant in $t$.
\end{proof}

Write $\mu_{s,d}=\kappa_{s,d}/(s-1)=m/b$ in lowest terms. We have $\mu_{s,d}\le d-1$, so \eqref{eq:endpoint-condition} holds with equality in its second inequality. Restrict to $h\equiv1\pmod b$. Theorem~\ref{thm:critical} and Lemma~\ref{lem:extraction} give
\begin{equation}\label{eq:critical-exponent}
R(s,t)\ge\Omega_{s,d}\bigl(t^{\mu_{s,d}/\log_2d}\bigr).
\end{equation}
For $d\ge s+1$, the codimension is constant and the denominator increases. The choice $d=s$ has the same codimension with a smaller denominator. Thus the optimized exponent is
\begin{equation}\label{eq:Ecrit}
E_{\rm crit}(s)=\max_{3\le d\le s}
\frac{(d-1)(2s-d)}{2(s-1)\log_2d}.
\end{equation}
Taking $d=s$ gives $s/(2\log_2s)$ with no rounding loss. For $d\le s+1$, the original marked condition permits only integers $m<\mu_{s,d}$. At a fixed $d$ its largest admissible label dimension is the largest such integer, when positive. The critical construction and restriction of scalars replace it by the exact rational endpoint.
% END SOURCE_BLOCK scalars

% BEGIN SOURCE_BLOCK frobenius_words
\section{Short Frobenius relations and the rank exponent}\label{sec:words}
The rank bound can be smaller than $d^H$ when point coordinates satisfy relations compatible with binary Frobenius. Let
\[
P=(1,X_1,\ldots,X_a,Y_1,\ldots,Y_r),\qquad d=1+a+r,
\]
where $a\ge1$, and put $X_0=1$. Suppose each $Y_i$ is a linear combination, over the flag field, of the functions
\begin{equation}\label{eq:short-relation}
\prod_{j=0}^{k-1}X_{\sigma(j)}^{2^j},\qquad \sigma:\{0,\ldots,k-1\}\longrightarrow\{0,\ldots,a\},\qquad k\ge2.
\end{equation}
The coefficients of these relations may be arbitrary. The assertion below is uniform over them.

\begin{lemma}[Word reduction]\label{lem:words}
If $r>0$, let $\rho=\rho(d,r,k)$ be the largest positive root of
\begin{equation}\label{eq:rho}
x^k-dx^{k-1}+r=0.
\end{equation}
Then $d-1<\rho<d$, and the graph's independence number is at most $d^k\rho^H$. If $r=0$, the bound is $d^H$ and we put $\rho=d$.
\end{lemma}

\begin{proof}
Represent a feature in $\Wcal_H$ by a word of length $H$ on the $d$ point-coordinate letters, with position $j$ contributing the $2^j$-th power of its letter. Whenever $Y_i$ is followed by $k-1$ constant letters, substitute \eqref{eq:short-relation}, raised to the power for that position. Each resulting word contains one fewer $Y$ letter. Repeated substitution terminates, leaving words that avoid every block $Y_i0^{k-1}$. Thus $\dim\Wcal_H$ is bounded by their number $u_H$.

No two forbidden blocks overlap: all letters after the initial $Y_i$ are constant letters. Appending a letter therefore gives
\[
u_n=d^n\quad(0\le n<k),\qquad u_n=du_{n-1}-ru_{n-k}\quad(n\ge k),
\]
and $\sum_{n\ge0}u_nz^n=(1-dz+rz^k)^{-1}$. One may also count with an automaton. Its states are a safe state and $k-1$ states recording a terminal $Y$ followed by $j$ zero letters, $0\le j\le k-2$. Any of the $a$ letters $X_i$ goes to the safe state; any of the $r$ letters $Y_i$ goes to the first active state. A zero stays safe or advances an active state, except that the last such transition is forbidden. The transition matrix is primitive, since every state reaches the safe state, every state is reachable from it, and the safe state has a loop. Its Perron root is the dominant growth root in \eqref{eq:rho}.

The polynomial in \eqref{eq:rho} is negative at $d-1$ and positive at $d$: here $r\le d-2<(d-1)^{k-1}$. Its derivative is positive for $x\ge d-1$, so it has exactly one root there. Every row sum of the transition matrix is either $d$ or $d-1$, and both occur. Irreducibility therefore places its Perron root strictly between these values. The unique root in that interval is the Perron root. For an explicit prefactor, let $v$ be a positive right Perron vector. Any two states are joined by a path of length at most $k$, each with positive integral weight. Hence $v_i/v_j\le\rho^k\le d^k$. Comparing the all-ones vector with $v$ gives $u_H\le d^k\rho^H$. Lemma~\ref{lem:rank} completes the proof.
\end{proof}

The compressed features in $a$ variables satisfy \eqref{eq:short-relation} when $k=\lceil\log_2s\rceil$. Indeed, for $j<s\le2^k$, expand $\ell_T(X)^j$ according to the binary digits of $j$. Every occupied digit contributes one coordinate raised to the corresponding binary power; unoccupied digits contribute $X_0$. Taking a coefficient of $T$ preserves this form. This gives short relations for polynomial outputs in $X$, although a direct interpolation family in the original labels need not have the required relation to its own point coordinates.

\subsection{A two-stage interpolation family}\label{sec:two-stage}
Fix $3\le d\le s$, write $\mu_{s,d}=m/b$ in lowest terms, and set
\begin{equation}\label{eq:two-stage-parameters}
a=\lceil\mu_{s,d}\rceil,\qquad r=d-1-a,\qquad k=\lceil\log_2s\rceil.
\end{equation}
First copy the $m$ base coordinates of the label into $m$ of the $ba$ base coordinates of $X\in E_q^a$. Parameterize the other $ba-m$ coordinates by independent compressed polynomials in the labels. Next parameterize all $r$ coordinates $Y_i$ and all $d-2$ slopes by independent $E_q$-valued compressed polynomials in $X$. The coefficient dimension is
\begin{equation}\label{eq:two-stage-L}
L=(ba-m)M(m,s)+b(r+d-2)M(a,s).
\end{equation}
The second-stage $Y_i$ obey \eqref{eq:short-relation}. The first-stage functions may have higher Frobenius complexity; the word count uses only the relation of $Y$ to $X$.

\begin{proposition}\label{prop:two-stage}
The universal ordered-clique incidence for this family, restricted to pairwise distinct $X$-tuples in each split geometric factor, has dimension at most
\[
L+ms-b\kappa_{s,d}\le L+m.
\]
It contains every actual ordered clique on distinct labels whenever $E_q$ is a field. Consequently there are $K_s$-free graphs on $q^{m-1}(q-\delta)$ vertices with
\[
\alpha\le C d^k\rho^{bh},
\]
for fixed $C,\delta$ and every sufficiently large $h\equiv1\pmod b$. They have deterministic preprocessing and polynomial-time local access for fixed $s,d$.
\end{proposition}

\begin{proof}
After algebraic closure, $E_q$ splits into $b$ factors. Restrict the configuration space to the open set where the $a$-tuples $X$ at different positions are distinct in every factor, as well as the open set of distinct base labels. Fix a full configuration there. The first-stage evaluation conditions have rank $s$ for each of their $ba-m$ outputs. In each split factor, the second-stage evaluation conditions also have rank $s$, by Lemma~\ref{lem:interpolation} applied to the distinct $X$-tuples. Once the full configuration is fixed, its $X$-values are fixed as well. The second-stage conditions are therefore linear in their own coefficients and do not impose additional nonlinear conditions on the first-stage coefficients. After splitting the algebra, changing from its base coordinates to its factor coordinates is an invertible linear transformation. The coefficient blocks in the two stages and in the split factors are independent. The solution fiber has dimension
\[
L-s\bigl((ba-m)+b(r+d-2)\bigr)=L-s\bigl(b(2d-3)-m\bigr).
\]
Cover both evaluation matrices by their nonzero-minor charts, as in Lemma~\ref{lem:universal-dimension}. The scalar configuration dimension is $s b(2d-3)-b\kappa_{s,d}$, giving the claimed bound.

The open restriction has a prime-field polynomial description. For every pair $i<j$, introduce $a$ algebra-valued witnesses and require
\[
\sum_{\ell=1}^a U_{ij,\ell}(X_{i,\ell}-X_{j,\ell})=1\quad\text{in }E_q.
\]
In each split factor this is solvable precisely when that difference vector is nonzero. Base-label distinctness has the analogous scalar witnesses. The witness variables serve only to describe the open set; their dimensions are not included in the incidence bound.

For an actual field $E_q$, distinct labels give distinct $X$-tuples because the copied base coordinates are retained. Every embedding into an algebraic closure is injective, so the difference remains nonzero in every split factor. Thus every actual ordered clique lies in the restricted incidence. Theorem~\ref{thm:diagonal-separation} deletes one of its edges, and Lemmas~\ref{lem:filter-rank} and \ref{lem:words} give the independence estimate. All polynomial families and filters are fixed before $h$ varies.
\end{proof}

\begin{corollary}\label{cor:frobenius-exponent}
With $\rho$ defined by \eqref{eq:two-stage-parameters} and \eqref{eq:rho}, put
\begin{equation}\label{eq:EF}
E_{\rm F}(s)=\max_{3\le d\le s}\frac{\mu_{s,d}}{\log_2\rho}.
\end{equation}
Then explicit witnessing graphs give $R(s,t)\ge\Omega_s(t^{E_{\rm F}(s)})$. In particular, $d=13$ for $s=16$ gives $\mu=38/5$, $a=8$, $r=4$, $k=4$, and
\begin{equation}\label{eq:R16}
R(16,t)\ge\Omega\left(t^{\frac{38}{5\log_2\rho_{16}}}\right),\qquad
\rho_{16}^4-13\rho_{16}^3+4=0,\quad 12<\rho_{16}<13.
\end{equation}
The exponent is $2.0539221767\ldots$.
\end{corollary}

\begin{proof}
Apply Lemma~\ref{lem:extraction} to Proposition~\ref{prop:two-stage}; $d^k$ is fixed. For the displayed instance, $\kappa_{16,13}=114$, so $m=38$, $b=5$ and $H=5h$, with $h\equiv1\pmod5$. The vertex count is of order $2^{38h}$ and the independence estimate is a fixed multiple of $\rho_{16}^{5h}$. Their ratio of logarithmic growth rates is the stated exponent.
\end{proof}

\begin{proposition}\label{prop:asymptotic}
As $s\to\infty$,
\[
E_{\rm crit}(s)=\frac{s}{2\log_2s}\left(1+\frac{1}{4(\ln s)^2}+O((\ln s)^{-3})\right),
\qquad E_{\rm F}(s)=\left(\frac12+o(1)\right)\frac{s}{\log_2s}.
\]
The maximum defining $E_{\rm crit}$ is attained at
$d=s-s/(2\ln s)+O(s/(\ln s)^2)$, up to integral rounding.
\end{proposition}

\begin{proof}
For the first assertion, treat $d$ as real in \eqref{eq:Ecrit}. Choices with $d\le s/(\ln s)^2$ give $o(s/\log s)$ by $\mu\le d-1$. Above that range, $\ln d=(1+o(1))\ln s$ uniformly. If also $d/s\le1-\varepsilon$, the normalized numerator loses at least a fixed proportion $\varepsilon^2$. Since $d=s$ gives $s/(2\log_2s)$, a maximizer has $d/s\to1$. Write $d=s(1-x)$ and $L_s=\ln s$. Apart from a relative $O(1/s)$ error, the ratio to $s/(2\log_2s)$ is
\[
\frac{1-x^2}{1+\ln(1-x)/L_s}.
\]
The stationary equation for the displayed approximation is $2x(L_s+\ln(1-x))=1+x$. As $x=o(1)$, it gives $x=1/(2L_s)+O(L_s^{-2})$. Expansion at this value gives $1+1/(4L_s^2)+O(L_s^{-3})$. The integral mesh has size $1/s$ and changes neither expansion.

For the second assertion, $\rho\le d$ gives $E_{\rm F}\ge E_{\rm crit}$. Also $\mu_{s,d}\le\min(d-1,s/2)$ and $\rho>d-1$ when $r>0$, while $\rho=d$ when $r=0$. If $d\le s/(\ln s)^2$, these bounds give $o(s/\log s)$, using that $x/\log_2x$ is increasing for $x$ sufficiently large and treating bounded $d$ separately. In the remaining range, $\log_2(d-1)=(1+o(1))\log_2s$, uniformly, and $\mu\le s/2$. This proves the upper bound.
\end{proof}
% END SOURCE_BLOCK frobenius_words

% BEGIN SOURCE_BLOCK finite_blocks
\subsection{Finite-block rank certificates}\label{sec:finite-blocks}
Short word relations give one usable upper bound on the feature dimension. A fixed-length symbolic rank computation gives another.

\begin{proposition}\label{prop:finite-block}
For a fixed point family, write $\beta_n=\dim\Wcal_n$. Then $\beta_{u+v}\le\beta_u\beta_v$. Consequently the limit $\eta=\lim_{n\to\infty}\beta_n^{1/n}=\inf_n\beta_n^{1/n}$ exists. More concretely, for every fixed $\ell\ge1$,
\[
\beta_H\le d^{\ell-1}\beta_\ell^{\lfloor H/\ell\rfloor}.
\]
If $\beta_\ell^{\rm gen}$ is the rank of the length-$\ell$ features over the rational function field of the coefficients, every specialization has rank at most $\beta_\ell^{\rm gen}$. In any compatible critical construction with $\beta_\ell^{\rm gen}>1$, this yields the Ramsey exponent $\mu\ell/\log_2\beta_\ell^{\rm gen}$.
\end{proposition}

\begin{proof}
Split a feature word into its first $u$ letters and last $v$ letters. It is a product of a member of $\Wcal_u$ and a $2^u$-th power of a member of $\Wcal_v$. Products of bases span a space of dimension at most $\beta_u\beta_v$. The field is perfect, so Frobenius does not change the dimension. Subadditivity of $\log\beta_n$ gives the limit and infimum; splitting into blocks of length $\ell$ and a remainder gives the finite inequality.

For a polynomial point family, expand the feature polynomials in the input monomial basis. The resulting finite matrix has polynomial entries in the coefficients. Every minor above its generic rank is identically zero and remains zero under specialization. Thus the generic rank is a uniform upper bound, including for exceptional coefficient choices. Since $\beta_\ell^{\rm gen}\ge1$, the finite-block inequality is bounded by a fixed multiple of $(\beta_\ell^{\rm gen})^{H/\ell}$. Apply Lemma~\ref{lem:extraction}.
\end{proof}

The limit by itself does not imply $\beta_H=O(\eta^H)$. The finite-block statement supplies the prefactor needed for an exact fixed-$s$ power of $t$. No additional numerical improvement from a symbolic block computation is asserted here.
% END SOURCE_BLOCK finite_blocks

% BEGIN SOURCE_BLOCK general_patterns
\section{Other forbidden configurations and fixed characteristic}\label{sec:general}
The critical argument depends on the configuration codimension, universal evaluation rank, and a distinguished edge. Flags supply these data, but the filtering step does not depend on their geometry.



\begin{proof}[Proof of Theorem~\ref{thm:configuration}]
Use compressed features through degree $\max_i v_i-1$. Universal interpolation, applied separately to each family, gives incidence dimension at most $L+mv_i-b\kappa_i\le L+m$. Relabel the positions of its distinguished edge as the remembered pair. The proof of Theorem~\ref{thm:diagonal-separation} produces $F_i$ vanishing on that pair image, with $F_i(A,Z,Z)\not\equiv0$.

The product $F=\prod_iF_i$ has nonzero diagonal restriction because the coefficient--label polynomial ring is a domain. Specialize a nonzero coefficient of that restriction and decode a fixed-degree polynomial complement. An occurrence of family $i$ loses its distinguished edge through the factor $F_i$. Lemma~\ref{lem:filter-rank} gives a fixed multiplier for the product. Finitely many orderings can be included as separate families. No connectedness assumption on a pattern is needed.
\end{proof}

\subsection{A fixed prime characteristic}
For a fixed prime $p$, the same interpolation, scalar restriction, diagonal separation and polynomial-complement algorithms work over $\F_p$. The degree of scalar restriction is again taken coprime to the varying extension degree. The rank denominator changes.

\begin{proposition}\label{prop:prime}
Put $B_{d,p}=\binom{d+p-2}{p-1}$. Over a flag field of order $p^H$, the assignment-uniform independence certificate has rank at most $B_{d,p}^H$. The critical construction therefore gives, for fixed $s,d,p$,
\[
R(s,t)\ge\Omega_{s,d,p}\left(t^{\kappa_{s,d}/((s-1)\log_p B_{d,p})}\right).
\]
\end{proposition}

\begin{proof}
Replace the binary Frobenius product by
\[
K^{p^H-1}=\prod_{j=0}^{H-1}\bigl(K^{p-1}\bigr)^{p^j}.
\]
Expanding $K^{p-1}$ uses the $B_{d,p}$ point monomials of degree $p-1$ in $d$ coordinates. Their multinomial coefficients are nonzero modulo $p$. The normalized constant coordinate absorbs the all-ones term exactly as in Lemma~\ref{lem:rank}. Thus there are at most $B_{d,p}^H$ features. The matrix $1-K^{p^H-1}$ is still triangular with diagonal $1$ on independent sets.

Prime-field descent uses $p$-th power Frobenius. To split the rational roots in the decoder, replace the binary trace tests by $\Tr_{\F_{p^h}/\F_p}(\beta_jX)-c$ for all $c\in\F_p$. The trace pairing is nondegenerate, and $p$ is fixed, so the cost remains polynomial in $h$. Apply the scalar critical construction with $\mu=m/b=\kappa_{s,d}/(s-1)$ and extract along $h\equiv1\pmod b$ as in Lemma~\ref{lem:extraction}, using logarithms to base $p$.
\end{proof}
% END SOURCE_BLOCK general_patterns

% BEGIN SOURCE_BLOCK closing
\section{Construction recipe and further questions}\label{sec:recipe}
For fixed $s,d$, the construction can be implemented in four stages. In preprocessing, reduce $\mu_{s,d}$ to $m/b$, choose the one- or two-stage coordinate family and a binary basis of $E_0$, and compute either a diagonal-separating filter or the finite-fiber polynomials. This stage terminates independently of $h$; no polynomial bound in $s,d$ is asserted. For an admissible $h$, initialize $\F_{2^h}$ and the fixed algebra $E_q$. Then specialize coefficients using a nonzero coefficient of the required deletion polynomial. Finally initialize its coordinate change and use the decoder and edge test, computing conflict colors locally if the quantified construction is chosen. The last three stages have deterministic bit cost polynomial in $h$ for fixed parameters.

The indexing domain is the product in \eqref{eq:product-labels}. Given two indices, decode their labels, compare them in the fixed order, evaluate the corresponding flags, and apply the incidence test together with the chosen filter or color rule. Neither enumerating the full vertex set nor finding a largest color class is part of a local query.

The main geometric limitation is now explicit. The unrestricted ordered flag configuration space has the exact codimension $\kappa_{s,d}$; refining the complete-flag count cannot improve it. A further gain must come from restricting configurations through a new coordinate family, finding a smaller rank growth rate compatible with universal interpolation, or proving diagonal separation beyond the dimension criterion. Properness of the pair image alone does not suffice for the last approach.

Finite-block ranks offer a concrete algebraic problem: compute $\beta_\ell^{\rm gen}$ for a compatible family and compare its growth with the short-word bound. A useful improvement would require a proved rank certificate and the same control of all distinct forbidden tuples. A low-rank coordinate family alone does not establish clique avoidance, because its incidence may acquire large exceptional strata.

\paragraph{Acknowledgements and declaration of assistance.}
The author thanks Sam Matheus for bringing this problem to his attention. Most of the ideas in this paper originated in exchanges with GPT-6 Pro. GPT-6.1 sol was used in Codex to assist with writing and revising the manuscript. The author has checked all arguments and computations and takes full responsibility for the contents of the paper.
% END SOURCE_BLOCK closing

% BEGIN SOURCE_BLOCK appendix_start
\appendix
\section{Dense indexing of a polynomial complement}\label{sec:decoder}
% END SOURCE_BLOCK appendix_start

% BEGIN SOURCE_BLOCK decoder
\begin{lemma}\label{lem:decoder}
Let $F\in\F_q[Z_1,\ldots,Z_m]$ be nonzero of degree at most $\Delta$, with $q=2^h>\Delta$. There is a deterministically decodable subset of $F\ne0$ of exactly $q^{m-1}(q-\Delta)$ points. For fixed $m,\Delta$, decoding and testing membership in this subset take time polynomial in $h$.
\end{lemma}

\begin{proof}
Put $\delta=\deg F$. If $\delta=0$, take the first $q-\Delta$ encoded field elements in each last-coordinate fiber. Suppose $\delta>0$ and let $F_\delta$ be its highest homogeneous part. Choose the first $\delta+1$ field elements as a grid in each coordinate. Some grid vector $u$ has $F_\delta(u)\ne0$: otherwise repeated application of the univariate root bound would make $F_\delta$ the zero polynomial. This search has size depending only on the fixed parameters. Since $u\ne0$, complete it deterministically to a basis by testing standard basis vectors in order. Let $S$ be the resulting invertible matrix with last column $u$.

Write $Z=S(y,x)^T$, with $y\in\F_q^{m-1}$. In $\widehat F(y,x)=F(S(y,x)^T)$ the coefficient of $x^\delta$ is the constant $F_\delta(u)\ne0$. Every fiber polynomial $f(X)=\widehat F(y,X)$ is therefore nonzero of degree $\delta$, with at most $\Delta$ rational roots.

Compute
\[
R(X)=\gcd(f(X),X^q-X).
\]
The polynomial $X^q-X$ has exactly the field elements as simple roots, so $R$ is the product of the distinct linear factors corresponding to rational roots of $f$. It is not necessary to form $X^q-X$: $h$ repeated squarings compute $X^q$ modulo $f$, which suffices for the gcd. If $R=1$, there are no roots to exclude.

Choose a basis $\beta_1,\ldots,\beta_h$ of $\F_q$ over $\F_2$. Split each current factor using gcds with
\begin{equation}\label{eq:trace-split}
T_j(X)=\sum_{\ell=0}^{h-1}(\beta_jX)^{2^\ell}
\quad\text{and}\quad T_j(X)-1,
\end{equation}
computed modulo that factor. On field elements $T_j$ is the trace and takes values $0,1$. Distinct roots are separated by some test because the trace pairing is nondegenerate. Indeed, $X+X^2+\cdots+X^{2^{h-1}}$ is a nonzero polynomial of degree less than $q$, so it has a value $w$ with nonzero trace. For any $v\ne0$, $\Tr((w/v)v)=\Tr(w)\ne0$. Linearity implies that some basis element pairs nontrivially with $v$. Apply this to the difference of two roots. After the $h$ tests, every remaining factor is linear. All degrees are at most $\Delta$, so root listing has polynomial bit cost in $h$ for fixed parameters.

Identify field elements with integers $0,\ldots,q-1$ by their binary polynomial-basis encoding. Use the product index set
\begin{equation}\label{eq:product-labels}
\F_q^{m-1}\times\{0,1,\ldots,q-\Delta-1\}.
\end{equation}
Given $(y,j)$, sort the excluded root indices in that fiber. Start at $x=j$ and pass through this list. If the next forbidden index is at most the current $x$, increase $x$ by one. Each increment compensates for one forbidden entry preceding the desired nonroot. At the end $x$ is the $j$th nonroot, counting from zero, and $x\le j+\Delta\le q-1$. Apply $S$ to $(y,x)^T$ to return to the original label. The map is injective, within and between fibers.

To test membership of a label, apply $S^{-1}$, reject a rational root, and subtract the number of smaller excluded indices from its last-coordinate index. The label belongs to the selected subset precisely when this nonroot index is less than $q-\Delta$. The same bounded root list computes this inverse. If $m=1$, the first factor of \eqref{eq:product-labels} is the one-element set of empty vectors. The prescription still applies.
\end{proof}

An efficient test for $F(z)\ne0$ alone would not provide this decoder. The coordinate change ensures a nonzero polynomial in every fiber, and the trace splitting lists only its bounded set of forbidden entries.
% END SOURCE_BLOCK decoder

% BEGIN SOURCE_BLOCK original_statements
\section{The original strict constructions}\label{sec:original-statements}

Set
\begin{align}\label{eq:original-parameters}
&s\ge3,\quad d\ge3,\quad 1\le r\le d-2,\quad m=d-1-r,\quad D_0=\binom d2,\notag\\
&c=r+d-2,\quad M=s+(m-1)\binom s2,\quad L=cM,\notag\\
&a_0=s(s+c-1),\quad \Delta_0=s^{a_0},\quad B_0=\Delta_0+1,\quad h_0=B_0^L.
\end{align}
Here $c$ is the number of polynomial outputs, $M$ the number of interpolation features per output, and $L$ the total number of coefficient slots. Features that vanish or become redundant in characteristic two are still counted as slots. The integer $a_0$ counts equations in a closed witness model, and $\Delta_0$ bounds its cumulative degree.

\begin{theorem}[Direct construction]\label{thm:direct}
If $rs>D_0$, then for every integer $h\ge h_0$ there is a deterministically specified graph on $(\F_{2^h})^m$ with
\[
|V|=2^{hm},\qquad \omega<s,\qquad \alpha\le d^h.
\]
The bounds hold for every total ordering in the adjacency rule. For fixed $s,d,r$, field initialization and adjacency take deterministic time polynomial in $h$. The coefficient prescription requires neither elimination nor a search for a suitable graph.
\end{theorem}

\begin{theorem}[Marked deletion]\label{thm:marked}
If $r(s-1)>\binom{d-1}2$, then for every $h\ge h_0$, with $q=2^h$, there is a deterministically specified graph with
\[
|V|=q^{m-1}(q-\Delta_0),\qquad \omega<s,\qquad \alpha\le d^h.
\]
Its labels form an explicitly indexed product set. For fixed $s,d,r$, decoding and adjacency take deterministic time polynomial in $h$. Finite algebraic preprocessing depends only on $s,d,r$; no efficient bound in those parameters is asserted.
\end{theorem}

The vertex count is positive, since $h_0>\Delta_0$ and $q>2\Delta_0$. The deletion argument first produces a polynomial complement; the decoder selects exactly $q-\Delta_0$ labels in every one of $q^{m-1}$ fibers. The normalization of one point coordinate absorbs the all-ones matrix into an existing rank feature, giving $d^h$ rather than $d^h+1$.

\begin{corollary}\label{cor:original-ramsey}
Under either theorem's hypotheses, explicit witnessing graphs give
\[
R(s,t)\ge\Omega_{s,d,r}\bigl(t^{(d-1-r)/\log_2 d}\bigr).
\]
Taking $d=s$ and $r=\lceil s/2\rceil$ in the direct construction gives exponent $(\lfloor s/2\rfloor-1)/\log_2 s$ for $s\ge4$. Taking $d=s$ and $r=\lfloor s/2\rfloor$ in the marked construction gives exponent $(\lceil s/2\rceil-1)/\log_2 s$ for $s\ge3$.
\end{corollary}

Here $f(t)\ge\Omega_{s,d,r}(g(t))$ means $f(t)\ge Cg(t)$ for every sufficiently large $t$, for some positive constant depending only on the displayed fixed parameters.
% END SOURCE_BLOCK original_statements

% BEGIN SOURCE_BLOCK original_direct
\subsection{Proofs}\label{sec:strict}
\subsubsection{A closed witness model}
Use the parameters \eqref{eq:original-parameters}. Take the $L$ coefficients and the $2d-3$ chart coordinates of each of the $s$ flags. At position $j$, let $z_j$ be the first $m$ point coordinates. Add the distinctness witnesses \eqref{eq:inverse-witness} for every pair. The closed set $W$ is defined by the following equations: $\binom s2$ clique incidences of degree at most two; $\binom s2$ distinctness equations of degree two; and $cs$ polynomial-graph equations equating output coordinates with their compressed feature combinations. The latter have total degree at most $s$, including coefficients. Thus there are
\[
2\binom s2+cs=s(s+c-1)=a_0
\]
equations, each of degree at most $s$, and $\cdeg W\le s^{a_0}=\Delta_0$.

Forgetting the witnesses maps $W$ onto the distinct-label coefficient--configuration incidence $\Ical$ on algebraically closed points. They have the same projected images in coefficients alone, or coefficients and the first label. The projection and descent assertions of Lemma~\ref{lem:degree} therefore give
\begin{align}\label{eq:exceptional-original}
&0\ne Q_A(A)\in\F_2[A],\quad\deg Q_A\le\Delta_0
&&\text{if }rs>\binom d2,\notag\\
&0\ne Q(A,Z)\in\F_2[A,Z],\quad\deg Q\le\Delta_0
&&\text{if }r(s-1)>\binom{d-1}2.
\end{align}
The first vanishes on every bad coefficient tuple, and the second on every bad tuple with its first label. Properness comes from \eqref{eq:original-dimension}; witness dimensions enter only the degree bound.

\subsection{Bounded-degree avoidance without elimination}
Choose an irreducible $f_h(T)\in\F_2[T]$ of degree $h\ge h_0$, and represent
\[
\F_q=\F_2[T]/(f_h(T)),\qquad\theta=T\pmod{f_h(T)}.
\]
The element $\theta$ has minimal-polynomial degree $h$. It need not generate the multiplicative group. In the fixed coefficient order prescribe
\begin{equation}\label{eq:kronecker}
A_j=\theta^{B_0^j}\qquad(0\le j<L).
\end{equation}

\begin{lemma}\label{lem:kronecker}
No nonzero polynomial in $\F_2[A_0,\ldots,A_{L-1}]$ of total degree at most $\Delta_0$ vanishes at \eqref{eq:kronecker}.
\end{lemma}

\begin{proof}
First substitute into a formal indeterminate:
\[
\widehat Q(T)=Q(T,T^{B_0},\ldots,T^{B_0^{L-1}}).
\]
A monomial $\prod_jA_j^{e_j}$ becomes $T^{\sum_je_jB_0^j}$. Since $\sum_je_j\le\Delta_0=B_0-1$, each $e_j$ is a base-$B_0$ digit. Distinct monomials give distinct exponents, so their coefficients cannot cancel and $\widehat Q\ne0$. Also
\[
\deg\widehat Q\le\Delta_0B_0^{L-1}<B_0^L=h_0\le h.
\]
The minimal-polynomial property gives $\widehat Q(\theta)\ne0$, which is the required value.
\end{proof}

\begin{proof}[Proof of Theorem~\ref{thm:direct}]
Use \eqref{eq:kronecker} and every label in $\F_q^m$. Copied coordinates recover each label, giving $q^m=2^{hm}$ vertices. Lemma~\ref{lem:rank} gives $\alpha\le d^h$. An ordered $s$-clique would put the coefficient tuple in the bad image and make $Q_A(A)=0$ in \eqref{eq:exceptional-original}, contrary to Lemma~\ref{lem:kronecker}. This excludes the algebraic ordered configurations independently of the chosen total order.

Generate $f_h$ by \cite{Shoup}, or use any supplied irreducible polynomial. Compute $A_0=\theta$ and $A_{j+1}=A_j^{B_0}$. For a label, form $\ell_T(z)$ over $\F_q$, compute its successive powers through $s-1$, and read the coefficient features. Evaluate the output combinations and \eqref{eq:kernel}. No expansion into all multivariate monomials is needed. With $s,d,r$ fixed, the number of operations and the bit length of $B_0$ are constant in $h$, and every field operation has polynomial bit cost. The exceptional polynomial proves correctness but is never computed by this algorithm.
\end{proof}

For definiteness use lexicographic order on binary polynomial-basis labels. The degree bound in \eqref{eq:exceptional-original} is independent of $h$, because the model is written over $\F_2$ before choosing the finite field. No equations $z^q-z=0$ enter the geometric argument.

\subsection{Marked deletion and preprocessing}
To compute $Q(A,Z)$, enumerate all nonzero prime-field polynomials of total degree at most $\Delta_0$, using a fixed monomial order and a fixed order of binary coefficient vectors. For each candidate append $vQ(A,z_1)-1=0$ to the witness equations for $W$. Accept the first candidate whose enlarged ideal contains $1$. The Nullstellensatz test described in Section~\ref{sec:degree-model} decides each candidate, and \eqref{eq:exceptional-original} guarantees acceptance. All this preprocessing depends only on $s,d,r$, with no efficient uniform bound asserted.

Write $Q(A,Z)=\sum_\gamma R_\gamma(A)Z^\gamma$. At least one $R_\gamma$ is nonzero, and each has degree at most $\Delta_0$. Under \eqref{eq:kronecker}, Lemma~\ref{lem:kronecker} keeps every nonzero coefficient nonzero. Hence
\[
F(Z)=Q(A,Z)
\]
is a nonzero polynomial of degree at most $\Delta_0$. Retain labels with $F(z)\ne0$. An ordered $s$-clique among them would have $Q(A,z_1)=0$, a contradiction. Restriction preserves the rank bound.

A nonzero degree-$\delta$ polynomial in $m$ variables has at most $\delta q^{m-1}$ zeros. For completeness, write it with last-variable degree $e$. Its nonzero leading coefficient has degree at most $\delta-e$ in the other variables. Induction bounds the choices making that coefficient zero by $(\delta-e)q^{m-2}$; allow all $q$ last coordinates there. On other fibers allow at most $e$ roots. The total is at most $(\delta-e)q^{m-1}+eq^{m-1}$. The case $m=1$ is the root bound, and a nonzero constant has no zeros. Thus deletion retains at least $q^m-\Delta_0q^{m-1}$ labels. Lemma~\ref{lem:decoder} selects exactly that many in a decodable product set. It proves the indexing and complexity assertions of Theorem~\ref{thm:marked}.

\begin{corollary}[Computed-polynomial specialization]\label{cor:adaptive-marked}
The conclusions of Theorem~\ref{thm:marked} hold already for
\[
h\ge\left\lceil\log_2(2\Delta_0+1)\right\rceil
\]
if the computed marked polynomial is specialized by Lemma~\ref{lem:adaptive}. The same replacement is available for the direct construction after computing $Q_A$, but then it uses algebraic preprocessing.
\end{corollary}

\begin{proof}
Choose a nonzero $R_\gamma$ and apply Lemma~\ref{lem:adaptive}. This keeps $F$ nonzero over a field with $q>2\Delta_0$. The preceding deletion and decoder arguments are unchanged. For the direct variant, apply the lemma to $Q_A$ itself.
\end{proof}
% END SOURCE_BLOCK original_direct

% BEGIN SOURCE_BLOCK original_extraction
\subsection{Original exponents and thresholds}\label{sec:original-exponents}
For Corollary~\ref{cor:original-ramsey}, choose $h=\lfloor\log_d(t-1)\rfloor$. If $t\ge d^{h_0}+1$, then $h\ge h_0$ and $\alpha\le d^h<t$. The direct graph has
\[
|V|=2^{hm}\ge2^{-m}(t-1)^{m/\log_2d}.
\]
The marked graph has at least half as many vertices, since $q>2\Delta_0$. This gives the claimed bound for every sufficiently large integer $t$.

For $d=s$, $r=\lceil s/2\rceil$ satisfies the direct condition when $s\ge4$, and $m=\lfloor s/2\rfloor-1$. In the marked case take $r=\lfloor s/2\rfloor$ and $m=\lceil s/2\rceil-1$. For $s=2u$, the difference $r(s-1)-\binom{s-1}2$ is $2u-1$; for $s=2u+1$ it is $u$. These are positive for the stated range.

Choosing the least admissible integer $r$ gives the optimized original exponents
\begin{align}\label{eq:old-optimized}
E_{\rm dir}(s)&=\max_d\frac{d-2-\lfloor d(d-1)/(2s)\rfloor}{\log_2d},\notag\\
E_{\rm mark}(s)&=\max_d\frac{d-2-\lfloor(d-1)(d-2)/(2(s-1))\rfloor}{\log_2d}.
\end{align}
Only $1\le r\le d-2$ enter these maxima. Admissibility implies $d<2s$ in the direct case and $d<2s-1$ in the marked case, so the ranges are finite. There is no admissible direct triple at $s=3$.

The marked inequality gives
\[
m<\frac{(d-1)(s-d/2)}{s-1}\le\frac s2.
\]
The last inequality is equivalent to $(d-s)(d-s-1)\ge0$, valid for integer $d$. Since $d\ge m+2$ and $x/\log_2(x+2)$ is increasing for $x>0$,
\[
\frac m{\log_2d}\le\frac{s/2}{\log_2(s/2+2)}.
\]
Together with $d=s$ this proves $E_{\rm mark}(s)=(1/2+o(1))s/\log_2s$. Every direct-admissible choice is marked-admissible, and the direct choice $d=s$ supplies the same leading asymptotic. This limits these sufficient-condition calculations, not all algebraic graph constructions.

For $d=s$ in the elimination-free direct prescription,
\[
M=O(s^3),\quad L=O(s^4),\quad a_0=O(s^2),\quad\log_2\Delta_0=O(s^2\log s).
\]
Thus $h_0=2^{O(s^6\log s)}$ and a sufficient Ramsey threshold is $2^{2^{O(s^6\log s)}}$. If $s\to\infty$ and $s^6\log s=o(\log\log t)$, the direct graph gives $R(s,t)\ge t^{(1/2-o(1))s/\log_2s}$. The improved range proved below uses computed polynomials and does not change this elimination-free prescription.
% END SOURCE_BLOCK original_extraction

% BEGIN SOURCE_BLOCK quantitative_fibers
\section{Quantitative availability through finite marked fibers}\label{sec:fibers}
Diagonal separation proves the endpoint without a prescribed separating degree. A different edge deletion gives explicit constants. It first bounds the number of forbidden tuples through each retained vertex and then colors their conflict graph. The polynomial used to retain vertices has a degree bound coming directly from the substituted clique equations.

\begin{lemma}[Finite marked fibers]\label{lem:finite-fibers}
Let $X\subseteq\A^L\times(\A^m)^s$ be a reduced prime-field algebraic set with
\[
\dim X\le L+m,\qquad\cdeg X\le D.
\]
There is a nonzero $Q(A,Z)\in\F_2[A,Z]$ of degree at most
\begin{equation}\label{eq:fiber-Delta}
\Delta=s m(s-1)D
\end{equation}
such that, when $Q(a,z)\ne0$, the fiber of $X$ fixing $A=a$ and any one marked label to $z$ has at most $D$ points. The polynomial is computable by finite deterministic preprocessing.
\end{lemma}

\begin{proof}
For every marked position $j$ and every scalar coordinate $Y$ in the other $s-1$ labels, project $X$ to $(A,Z_j,Y)$. Its closure is proper in the $(L+m+1)$-dimensional target and has cumulative degree at most $D$. Lemma~\ref{lem:degree} supplies a nonzero polynomial
\[
G_{j,Y}(A,Z,Y)=\sum_{u=0}^{u_{j,Y}}g_{j,Y,u}(A,Z)Y^u
\]
of degree at most $D$ vanishing on it. Choose its highest nonzero coefficient in $Y$ and take the product of those coefficients over all $s m(s-1)$ choices of $(j,Y)$. The polynomial ring is a domain, and each selected leading coefficient is nonzero, so their product is nonzero. A nonvanishing product at $(a,z)$ keeps every specialized coordinate equation nonzero simultaneously, for every marked position. Its degree is at most \eqref{eq:fiber-Delta}.

If $Q(a,z)\ne0$, each unmarked scalar coordinate of a point in the fiber is among finitely many roots of the corresponding nonzero univariate polynomial. When its degree is zero the fiber is empty. Otherwise the fiber lies in a finite Cartesian product and is zero-dimensional. It is a linear section of $X$, so its cumulative degree is at most $D$. A reduced zero-dimensional set has at most its degree many geometric points. These observations give the bound, including for nonreduced scheme-theoretic fibers.

Compute each $G_{j,Y}$ by enumerating binary polynomials of degree at most $D$ and testing vanishing on the incidence model. The degree and descent statements ensure termination. The leading coefficients and their product are then ordinary polynomial computations.
\end{proof}

Specialize $Q$ by Lemma~\ref{lem:adaptive} and use Lemma~\ref{lem:decoder}, with $q>2\Delta$, to obtain an indexed set $S$ of exactly $q^{m-1}(q-\Delta)$ labels. Form the auxiliary conflict graph on $S$ by joining every pair of vertices occurring together in an ordered $s$-clique of the original graph. Each fixed vertex can occur in any of $s$ marked positions, with at most $D$ tuples for each position. Its conflict degree is therefore at most
\begin{equation}\label{eq:Lambda}
\Lambda=s(s-1)D.
\end{equation}
Its neighbors can be listed locally. For each position, find the base-field roots of all the $G_{j,Y}(a,z,Y)$, enumerate their fixed-size Cartesian product, and test the clique equations, the fixed increasing order of the tuple, distinctness and membership in $S$. Discard tuples failing any test and remove repeated neighbors from the resulting list. Root finding and the membership test are supplied by the decoder algorithms. The enumeration size is a parameter-dependent constant, and the bit cost is polynomial in $h$ for fixed parameters.

\begin{lemma}[Local color reduction]\label{lem:color}
A graph of maximum degree at most $\Lambda$ on labels in $\F_{2^h}^m$, with a polynomial-time neighbor oracle for fixed $\Lambda,m$, has a deterministic locally computable proper coloring with at most
\begin{equation}\label{eq:palette}
C=(16(\Lambda+1))^{16}
\end{equation}
colors. The bit cost of computing a color is polynomial in $h$ for fixed $\Lambda,m$.
\end{lemma}

\begin{proof}
Start with the injective binary encoding of the labels, using the common palette bound $N_0=2^{mh}$. Suppose a proper coloring has palette bound $N\ge2$, and put $\ell=\lceil\log_2N\rceil$. Choose the smallest power of two $p>\Lambda\ell$ (or $p=2$ if necessary), so $p\le2(\Lambda\ell+1)$. Encode the old color as a polynomial $f$ over $\F_p$ of degree at most $\ell$, using its base-$p$ expansion. Distinct colors have distinct polynomials. At a vertex choose the first field element $x$ such that $f(x)$ differs from the value at $x$ of every neighbor's polynomial. The union of their root sets has size at most $\Lambda\ell<p$, so such an $x$ exists. Assign the new color $(x,f(x))$.

Two adjacent vertices assigned the same new color would have the same $x$ and equal values there, contrary to the choice at either endpoint. Thus the coloring is proper, and its palette bound is $p^2$. For $N\ge2$,
\[
p^2\le16(\Lambda+1)^2(\log_2N)^2.
\]
For $N\ge C$ this is at most $\sqrt N$. To check it at $C$, write $A=\Lambda+1\ge1$: the left side is $4096A^2(4+\log_2A)^2$, whereas $\sqrt C=(16A)^8$; the inequality follows already from $4+\log_2A\le4A$. The ratio $(\log_2N)^2/\sqrt N$ decreases throughout this range. Iterate until the palette bound is at most $C$. There are $O(\log h)$ rounds for fixed $m,\Lambda$.

Each round uses only the previous colors of a vertex and its neighbors. Recursively computing these colors explores at most $(\Lambda+1)^{O(\log h)}$ oracle calls, hence polynomially many for fixed $\Lambda$. The auxiliary fields have size polynomial in $h$; deterministic field initialization and their enumeration have the same cost. Each pair $(x,f(x))$ is encoded as an integer in the common palette of size $p^2$. Every vertex uses the same successive palette bounds and stops after the same number of rounds, determined from $N_0$, rather than from its own attained color. All rounding and ordering choices are fixed.
\end{proof}

Retain an original edge only when its endpoints have the same conflict color. An $s$-clique would be monochromatic, whereas any two of its vertices are adjacent in the properly colored conflict graph. Thus none remains. On an independent set, the vertices in each color class are independent in the original graph, so the independence estimate is multiplied by at most $C$. All color classes are retained; no global search for a largest class is required.

\begin{theorem}[Quantified critical construction]\label{thm:quantified}
For a critical interpolating flag family satisfying \eqref{eq:endpoint-condition}, let $e$ be the degree bound in Proposition~\ref{prop:substituted-degree} and put
\begin{equation}\label{eq:quantitative-constants}
D=(2e)^{b\binom s2},\qquad
\Delta=s m(s-1)D,\qquad\Lambda=s(s-1)D,\qquad C=(16(\Lambda+1))^{16}.
\end{equation}
For every $h$ with $\gcd(b,h)=1$ and $2^h>2\Delta$, finite deterministic preprocessing produces $K_s$-free graphs on $q^{m-1}(q-\Delta)$ vertices. Their independence number is at most $C d^{bh}$, or $C d^k\rho^{bh}$ for the two-stage family. Local access takes time polynomial in $h$ for fixed parameters.
\end{theorem}

\begin{proof}
Let $X$ be the closure of the constructible clique incidence in coefficient--label space, without inverse witnesses. For the two-stage family take the distinct-$X$ open restriction used in Proposition~\ref{prop:two-stage}; its closure contains every actual clique when $E_q$ is a field. The interpolation dimension bound gives $\dim X\le L+m$, and Proposition~\ref{prop:substituted-degree} gives $\cdeg X\le D$. Apply Lemma~\ref{lem:finite-fibers}, specialize its polynomial, decode $S$, and color the conflict graph by Lemma~\ref{lem:color}. The preceding edge rule excludes all $s$-cliques and has the stated independence multiplier. It uses only local polynomial computations, root finding, membership and the color oracle.
\end{proof}

\begin{corollary}\label{cor:growing-range}
There are deterministic witnessing graphs yielding
\[
R(s,t)\ge t^{(1/2-o(1))s/\log_2s}
\]
when $s\to\infty$ and $s^2(\log_2s)^2=o(\log_2t)$. For fixed $s$, a sufficient availability threshold in this construction has size $2^{O(s^2(\log s)^2)}$. The preprocessing and local-access constants still depend on $s$.
\end{corollary}

\begin{proof}
Use the single-stage family with $d=s$ and $\mu=s/2$. In lowest terms this has $b=1,m=s/2$ for even $s$, and $b=2,m=s$ for odd $s$. Take $e=s$. Then $\log_2D$, $\log_2\Delta$ and $\log_2C$ are $O(s^2\log s)$, and an admissible initial degree is $O(s^2\log s)$. Choose the largest $h\equiv1\pmod b$ with $Cd^{bh}\le t-1$. Once it is in the available range, the gap and vertex-density bounds give
\[
\log_2|V|\ge\frac{s}{2\log_2s}\bigl(\log_2(t-1)-\log_2C\bigr)-mb-1.
\]
The initial-degree requirement is met for $t\ge2^{O(s^2(\log s)^2)}$. Under the displayed little-oh hypothesis the losses are negligible relative to $(s/(2\log_2s))\log_2t$. This proves the assertion without asserting polynomial cost uniformly in $s$.
\end{proof}
% END SOURCE_BLOCK quantitative_fibers

% BEGIN SOURCE_BLOCK geometry_appendix_heading
\section{Geometric conventions and the complete-flag count}
% END SOURCE_BLOCK geometry_appendix_heading

% BEGIN SOURCE_BLOCK geometric_language
\subsection{Dimension conventions}
All geometric dimensions below are over $\overline{\F}_2$. An affine algebraic set is a common zero set of finitely many polynomials. These are the closed sets of the Zariski topology. Inequalities such as $f\ne0$ are open conditions. A locally closed set is a closed set restricted by an open condition; a constructible set is a finite union of locally closed sets. Incidence and exact-rank conditions are of this form, using equations and nonvanishing of minors.

For an irreducible affine set, its dimension is the transcendence degree of its coordinate field. For a reducible set, dimension is the maximum over its irreducible components. For a constructible set, dimension means that of its Zariski closure. Dimensions add under products, finite unions take the maximum dimension, and the closure of the image of a polynomial map cannot have larger dimension than its source. Parameter counts are thus valid when supported by finite algebraic chart covers. We use constructibility of images and the fiber-dimension theorem in their standard forms \cite[Chapter I]{Hartshorne}.

An image need not be closed: projecting $xy=1$ to the $x$-coordinate gives $\Bbbk\setminus\{0\}$, whose closure is $\A^1$. Nevertheless, an image of dimension less than the target dimension has proper closure. A generic choice means a choice in a specified nonempty Zariski open set; it does not mean a random choice.
% END SOURCE_BLOCK geometric_language

% BEGIN SOURCE_BLOCK complete_flag_count
\begin{lemma}\label{lem:complete-flag-bound}
The ordered configuration space satisfies
\[
\dim\Ccal_{s,d}\le\binom d2+s(d-2).
\]
\end{lemma}

\begin{proof}
Put $R_j=\spanop(P_1,\ldots,P_j)$. Ordered incidence and self-incidence imply $R_j\subseteq H_j$, so every $R_j$ is nonzero and proper. The distinct subspaces in the nested sequence of prefix spans extend to a complete flag
\[
F_1\subset F_2\subset\cdots\subset F_{d-1},\qquad\dim F_u=u.
\]
Complete flags have dimension $(d-1)+(d-2)+\cdots+1=\binom d2$, by successive projective-space bundles. Fix one and indices $u_j$ such that $R_j=F_{u_j}$. Allow independently all pairs with $P_j\subseteq F_{u_j}\subseteq H_j$. Points in $F_{u_j}$ and hyperplanes containing it form
\[
\PP^{u_j-1}\times\PP^{d-u_j-1},
\]
of dimension $d-2$. Allowing all $s$ choices and the complete flag gives dimension at most $\binom d2+s(d-2)$. Bases adapted to the flag supply local product charts. Finitely many index sequences cover the configurations, and images do not increase dimension.
\end{proof}
% END SOURCE_BLOCK complete_flag_count

% BEGIN SOURCE_BLOCK bibliography
\begin{thebibliography}{99}
\bibitem{AP} N. Alon and P. Pudl\'ak, Constructive lower bounds for off-diagonal Ramsey numbers, \emph{Israel J. Math.} \textbf{122} (2001), 243--251.
\bibitem{Bradac} D. Brada\v{c}, Off-diagonal Ramsey numbers, preprint (2026), \href{https://arxiv.org/abs/2605.28793v3}{arXiv:2605.28793v3}.
\bibitem{DKL} Z. Dvir, J. Koll\'ar and S. Lovett, Variety evasive sets, \emph{Comput. Complexity} \textbf{23} (2014), 509--529.
\bibitem{HIV} J. Han, F. Ihringer and H. Van Maldeghem, Optimally pseudorandom $K_4$-free graphs, preprint (2026), \href{https://arxiv.org/abs/2609.32461}{arXiv:2609.32461}.
\bibitem{Hartshorne} R. Hartshorne, \emph{Algebraic Geometry}, Graduate Texts in Mathematics 52, Springer, New York, 1977.
\bibitem{IM} F. Ihringer and S. Mattheus, An improved algebraic construction for Ramsey numbers, preprint (2026), \href{https://arxiv.org/abs/2608.21769}{arXiv:2608.21769}.
\bibitem{KPR} A. Kostochka, P. Pudl\'ak and V. R\"odl, Some constructive bounds on Ramsey numbers, \emph{J. Combin. Theory Ser. B} \textbf{100} (2010), 439--445.
\bibitem{LR} G. Lachaud and R. Rolland, On the number of points of algebraic sets over finite fields, \emph{J. Pure Appl. Algebra} \textbf{219} (2015), 5117--5136.
\bibitem{LNZ} J. Lim, J. Nie and J. Zeng, Evasive sets, twisted varieties, and container-clique trees, preprint (2025), \href{https://arxiv.org/abs/2507.07594}{arXiv:2507.07594}.
\bibitem{Shoup} V. Shoup, New algorithms for finding irreducible polynomials over finite fields, \emph{Math. Comp.} \textbf{54} (1990), 435--447.
\end{thebibliography}
\end{document}
% END SOURCE_BLOCK bibliography
